\section{Introduction}
Let $M\subset\R^k$ be a known residual image and $R>0$ an inverse information matrix. Consider
\begin{equation}\label{eq:distance}
 C_R(b)=\inf_{x\in M}(x-b)^{\mathsf T}R^{-1}(x-b).
\end{equation}
At a smooth point, its full second jet sees only the normal compression of $R$. A statistical model may constrain $R$ enough to remove this ambiguity. The inverse problem is therefore governed both by the geometry of $M$ and by the linear space of admissible information matrices.

We study the calibrated family
\[
 R_\theta=JH(\theta)J^{\mathsf T},\qquad
 H(\theta)_{ab}=\theta_{a+b},\quad 0\le a,b\le k-1,
\]
where $J$ is fixed and invertible, and $\theta$ lies in an open positive moment cone of dimension $p=2k-1$. Set $n=k-1$ and $V_n=H^0(\Pj^1,\mathcal O(n))$. Dehomogenization identifies real coefficient vectors with $\PP_n$. A normal space $N$ corresponds to $U=J^{\mathsf T}N$. The observations specify the moment functional on $U^2=\Span\{fg:f,g\in U\}$. Multiple contacts specify a sum of such spaces, without cross-contact products.

The central geometric issue is not only whether a general product map is onto, but how large its entire failure locus is. Put $q=\binom{c+1}{2}$ and
\[
 X_L=\Gr(c,V_n)^L,\qquad
 \Delta^{(L)}=\{(U_1,\ldots,U_L):\textstyle\sum_\nu U_\nu^2\ne V_{2n}\}.
\]
All codimensions in the following theorem are complex codimensions of reduced closed sets.
\begin{theorem}[Global failure and independent contacts]\label{thm:global-intro}
Suppose $4\le c<k$, $L\ge1$, and $Lq\ge p$. Then
\begin{equation}\label{eq:codimension}
 \codim_{X_L}\Delta^{(L)}=\min\{Lq-p+1,Lc-3\}.
\end{equation}
There is an irreducible common-secant subvariety $\Sigma_L\subset\Delta^{(L)}$ of codimension $Lc-3$. If $Lc-3\le Lq-p+1$, it is a maximal-dimensional irreducible component of $\Delta^{(L)}$. Consequently a general $L$-tuple is informative if and only if $Lq\ge p$.
\end{theorem}
The common-secant condition means that all $U_\nu$ lie in a hyperplane represented by a point of the secant variety of the evaluation rational normal curve in $\Pj(V_n^*)$. Thus the second term of \eqref{eq:codimension} is a concrete geometric obstruction, not a formal expected codimension. For one contact, the common-base-point locus lies properly inside this larger irreducible family and is not itself a component of the failure locus.

For the synchronized residual image
\[
 q_A(\eta)=((a_i^{\mathsf T}\eta)^2)_{i=1}^k,\qquad A\in\R^{k\times d},
\]
the normal dimension is $c=k-d$. Native realizability is an open positive-vector condition. Global separation at a prescribed normal space requires, in addition, that the corresponding tangent plane not split across a nontrivial coordinate partition. This exact criterion is proved in Theorem~\ref{thm:fixed-realization}.
\begin{theorem}[Sharp native single-contact threshold]\label{thm:threshold}
Fix $d\ge2$ and any separated root and clock configuration of Section~\ref{sec:native}. For $k>d$, there exist full-rank loadings $A$ with $q_A(e_1)>0$ for which the single full second jet $D^2C_{R_\theta}(q_A(e_1))$ identifies every $\theta\in\Omega$ if and only if
\[
 \binom{k-d+1}{2}\ge2k-1.
\]
In the sufficient range the assertion holds on a nonempty Zariski-open loading set, for the global distance function. In particular
\begin{equation}\label{eq:threshold}
 k_{\min}(d)=d+\left\lceil\frac{3+\sqrt{16d+1}}2\right\rceil
            =d+2\sqrt d+O(1).
\end{equation}
Unrestricted positive matrices at the same contact retain an invisible space of dimension $d(2k-d+1)/2$.
\end{theorem}
A contact here is a matrix-valued observation with $q$ entries, not a scalar datum. The comparison in \eqref{eq:threshold} is across calibrated families with varying $k$, fixed $d$, and arbitrary separated configurations at each $k$; it is not a universal threshold for every information model. Independently designed contacts can use different loadings with the same $J$. Several points of one fixed loading are not asserted to vary independently.

\subsection*{Method and relation to earlier theory}
Hankel duality with polynomial multiplication is classical \cite{MP}. The saturated one-space maximal-rank statement is also a consequence of classical postulation of general rational curves \cite{BE85,BE87,Larson}. Our proof of Theorem~\ref{thm:global-intro} instead combines the classical dimensions of Hankel rank strata \cite[Sections 1--2]{CMSV} with an isotropic-incidence count, an optimization over all ranks, and an explicit common-secant family. This proves the full codimension formula and the coupled multi-space assertion, not just generic surjectivity. The classical maximal-minor height bound supplies the other matching upper bound \cite[Introduction]{EHU}. These classical inputs are distinguished from the incidence optimization and native fixed-space realization proved here.

Polynomial product spaces also occur in the study of Gram spectrahedra \cite{Sch}; the present observation is the dual moment functional on within-contact products. The monomial specialization is a restricted additive-basis problem \cite{Yu,Kohonen14,Kohonen17}, and can have a strictly worse threshold. Quadratic sensing and phase retrieval concern reconstruction of the input to a quadratic measurement map \cite{BCE}; our unknown is instead an information metric, and injectivity on all of $\Sym_d$ is not assumed. Complete-conormal projective duality \cite{GKZ} addresses a different, larger observation category.

The statistical statements likewise keep their experiments separate. The coalesced-root model has a smooth positive extension, but its physical local parameter set is a cone. Its efficient Gaussian quotient and $N^{-1/4}$ root-amplitude scale are derived explicitly in Section~\ref{sec:boundary}. For observed Poisson exposures we form normal-compression score statistics directly from the counts. Their local information loss is exactly the loss of rank of the product map. This is a local likelihood statement in a specified target subexperiment, not a claim of global equivalence of all raw data and finitely many distance readings. The finite-sample estimators and computed finite-offset features in the appendices are retained as distinct procedures. Standard constrained and regular asymptotic principles are attributed to \cite{Geyer,vdV}.

\section{Exact fibres of prescribed contact data}\label{sec:contact}
For $\theta\in\R^p$ write
\[
 \ell_\theta\left(\sum_{a=0}^{2n}f_at^a\right)=\sum_a f_a\theta_a.
\]
Let $\Omega\subset\R^p$ be open with $R_\theta>0$. At each prescribed point $x_\nu$, assume that the global image has one embedded smooth germ and that its local minimizing projection computes \eqref{eq:distance}. Parameter differentiation below uses analytic germs and a common projection tube on compact metric sets. These conditions are verified for the native loading family in Section~\ref{sec:realization}.

Let $Z_\nu$ have orthonormal columns spanning the Euclidean normal space. Put
\[
 B_\nu=Z_\nu^{\mathsf T}R_\theta Z_\nu,\qquad
 G_\nu=\tfrac12Z_\nu^{\mathsf T}D^2C_{R_\theta}(x_\nu)Z_\nu.
\]
\begin{lemma}[Normal block]\label{lem:normal}
One has $G_\nu=B_\nu^{-1}$. The full Hessian has tangent kernel and is determined by this normal block.
\end{lemma}
\begin{proof}
The quadratic contact with a manifold equals that with its tangent space $T$. At the minimizer of $(u-v)^{\mathsf T}R^{-1}(u-v)$ over $v\in T$, the vector $R^{-1}(u-v)$ is normal. Hence $u-v=RZn$, and normal projection gives $Bn=Z^{\mathsf T}u$. The minimum is $(Z^{\mathsf T}u)^{\mathsf T}B^{-1}Z^{\mathsf T}u$. A local parametrization and its positive tangent Gram matrix justify the minimizing projection by the implicit function theorem and yield this quadratic term.
\end{proof}
Set $U_\nu=J^{\mathsf T}\ran Z_\nu$ in polynomial coordinates, and define
\[
 \mathcal U=\operatorname{im}\left(\bigoplus_{\nu=1}^L\Sym^2U_\nu
                  \xrightarrow{\ \sum\mathrm{multiplication}\ }\PP_{2n}\right),
 \qquad \LL\theta=(B_\nu(\theta))_\nu.
\]
The codomain of $\LL$ has the direct-sum Frobenius norm.
\begin{theorem}[Exact contact fibre]\label{thm:fibre}
The exact second-jet fibre through $\theta$ is
\[
 (\theta+\mathcal U^\perp)\cap\Omega.
\]
Its local dimension is $p-\dim\mathcal U$. Recovery on $\Omega$ is equivalent to $\mathcal U=\PP_{2n}$, or to injectivity of $\LL$.
\end{theorem}
\begin{proof}
For normal vectors $v,w$ with coefficient polynomials $f_v,f_w$,
\[
 v^{\mathsf T}(R_{\theta'}-R_\theta)w
       =\ell_{\theta'-\theta}(f_vf_w).
\]
Equality of the jets is equality of the positive normal blocks, by Lemma~\ref{lem:normal}. It is therefore precisely annihilation of $\mathcal U$. Inversion is a diffeomorphism on positive blocks; openness of $\Omega$ gives the asserted fibre dimension.
\end{proof}
For one $d$-dimensional germ there are at most $q=\binom{k-d+1}{2}$ independent products. For unrestricted matrices, fixing the normal--normal block leaves tangent--tangent and mixed blocks free, with total dimension
\begin{equation}\label{eq:ambient}
 \frac{k(k+1)-(k-d)(k-d+1)}2=\frac{d(2k-d+1)}2.
\end{equation}
Small perturbations in every such direction preserve positivity.

\section{The native information family}\label{sec:native}
Fix distinct strictly positive rank-one four-cell probability tables $M_1,M_2$, roots $0<r_1<\cdots<r_k<D$, and $m=2k+1$ clocks $D<T_1<\cdots<T_m$. Use
\[
 f_1(T)=f_{1,0}(T)\prod_i((T-r_i-s_i/2)^2-v_i),\qquad
 p_j=\pi_jM_1+(1-\pi_j)M_2,
\]
\[
 \frac{\pi_j}{1-\pi_j}=\frac{\alpha f_1(T_j)}{(1-\alpha)f_2(T_j)}.
\]
The fixed factors are positive at the clocks. At $s=v=0$, the free coordinates are log mixing odds and the root-sum displacements; the physical retained coordinates satisfy $v_i\ge0$. The likelihood is positive and analytic on a two-sided neighbourhood of this point. Its restriction to $v\ge0$, not that extension's unrestricted estimator, is the physical experiment.

Let $\lambda_j>0$ be information weights and set
\[
 \kappa_j=\pi_j^2(1-\pi_j)^2\sum_{a=1}^4
      \frac{(M_{1a}-M_{2a})^2}{p_{ja}},\qquad
 V_j=\big(1,-((T_j-r_i)^{-1})_i,-((T_j-r_i)^{-2})_i\big).
\]
The $k$ simple poles precede the $k$ double poles. Write
\[
 h(t)=\prod_i(t-r_i)^2,\quad p_*(t)=\prod_i(t-r_i),\quad P(t)=\prod_j(t-T_j),
\]
\begin{equation}\label{eq:nativeconstants}
 \begin{split}
 d_i&=\frac{P(r_i)}{\prod_{a\ne i}(r_i-r_a)^2},\qquad
 \xi_j=\frac{h(T_j)}{P'(T_j)},\\
 J_{ia}&=[t^a]d_i\frac{p_*(t)}{t-r_i},\qquad
 \rho_j=\frac{\xi_j^2}{p_*(T_j)^2}.
 \end{split}
\end{equation}
\begin{proposition}[Profiled quadratic information]\label{prop:native}
The retained block of the inverse full information, equivalently the inverse nuisance Schur complement, is
\begin{equation}\label{eq:native}
 R=JH(\theta)J^{\mathsf T},\qquad
 \theta=\sum_{j=1}^m\frac{\rho_j}{\lambda_j\kappa_j}w_{2n}(T_j),
 \quad w_s(t)=(1,t,\ldots,t^s)^{\mathsf T}.
\end{equation}
The matrix $J$ is invertible. Independent positive information weights give an open moment cone $\Omega$ of dimension $p=2k-1$.
\end{proposition}
\begin{proof}
Multiplying a linear combination of the columns of $V$ by $h$ gives a polynomial of degree at most $2k$. Vanishing at the $2k+1$ clocks forces it to be zero; principal parts at the roots and the constant at infinity then force every coefficient to vanish. Thus $V$ is invertible and the full information is $V^{\mathsf T}\diag(\lambda_j\kappa_j)V$.

For arbitrary log-odds increments $y_j$, interpolate
\[
 g(t)=\sum_jh(T_j)y_j\frac{P(t)}{(t-T_j)P'(T_j)}.
\]
The retained coefficient at $r_i$ is $-g(r_i)/\prod_{a\ne i}(r_i-r_a)^2$. Thus the retained rows $Z_0$ of $V^{-1}$ are
\[
 (Z_0)_{ij}=\frac{\xi_jd_i}{T_j-r_i}
        =\frac{\xi_j}{p_*(T_j)}(Jw_n(T_j))_i.
\]
Their weighted Gram matrix gives \eqref{eq:native}. Evaluation of the polynomials defining $J$ at the distinct roots is diagonal with nonzero entries, proving invertibility. The positive moment weights range over an open orthant. The moment matrix has rank $p$ by Vandermonde independence, so its image is open. Positive atomic weights at distinct clocks give positive definite $H(\theta)$.
\end{proof}
The weights in this proposition need not be unknown likelihood parameters. For a realized fixed allocation they are $n_j/N$; for a limiting allocation they are its limits; a different prespecified $\lambda$ defines a target design metric. Section~\ref{sec:score-contact} uses a further protocol in which exposures are unknown but their counts are observed. Section~\ref{sec:boundary} identifies \eqref{eq:native} as the covariance of the efficient Gaussian quotient of the physical boundary experiment.
